\scanpage{106}%
i.e. que la filtration dudit foncteur admet un\deleted{e} splittage compatible avec les produits tensoriels. Mais le choix d'un tel splittage équivaut à celui d'un point \deleted{sous} sur un certain torseur à droite $Q=Q'$ sous le \inserted{schéma en} groupes \inserted{(pro-unipotent)} $H$ des automorphismes de $F$ respectant la filtration et induisant l'identité sur le gradué associé; il n'est pas du tout canonique !

3) Appelons \uline{pré-structure de Hodge} sur un espace vectoriel \inserted{$V$} de dimension finie sur $\underline{Q}$, la donnée d'une bigraduation sur $V\unsure{\otimes|\Omega}_{\underline{Q}}\underline{C}=V_{\underline{C}}$, $V_{\underline{C}}=\handwritten{\coprod_{p,q}}V^{p,q}$, telle que a) la graduation totale correspondante soit ``définie sur $\underline{Q}$'' \uline{i.e.} $\handwritten{\coprod}_{p+q=n}V^{p,q}$ provienne d'un sous espace $V_{\underline{Q}}^n$ de $V$, et b) on a $\overset{\handwritten{\rule{.65em}{.35pt}}}{V}^{p,q}=V^{q,p}$, où $x\handwritten{\mapsto}\overset{\handwritten{\rule{.55em}{.35pt}}}{x}$ désigne la conjugaison complexe. \unsure{(}NB Généralisation à des corps plus généraux laissée à Saavedra). Les vectoriels $V$ munis d'une pré-structure de Hodge forment une catégorie tensorielle sur $\underline{Q}$ dans un sens évident, muni d'un foncteur fibre canonique, le foncteur ``oubli''\inserted{$F$}. On trouve donc un groupe de Galois $G$, et \deleted{tout} plus généralement toute \inserted{$\otimes$-}sous-catégorie $\underline{C}$ de la catégorie précédente nous définit un groupe de Galois $G$. La bigraduation sur le foncteur \deleted{oubli} $F_{\underline{C}}(V)=V\otimes_{\underline{Q}}\underline{C}$ correspond, en vertu de 1) (où il convient cependant de se permettre une extension du corps de base sur le foncteur fibre envisagé ....) d'un homomorphisme ${\underline{G}_m}_{\underline{C}}^2\handwritten{\longrightarrow}G_{\underline{C}}$, i.e. de deux homomorphismes \inserted{qui commutent}
\[
i_1,i_2:{\underline{G}_m}_{\underline{C}}\handwritten{\longrightarrow}G_{\underline{C}}.
\]
Les deux conditions a) et b) imposées aux structures envisagées s'interprètent respectivement par les faits que l'homomorphisme
\[
i_{\underline{C}}=i_1i_2\unsure{\doteqdot|\div|=}\handwritten{\lambda\mapsto}i_1(\handwritten{\lambda})i_2(\handwritten{\lambda})
\]
est ``défini sur $\underline{Q}$'' i.e. provient d'un homomorphisme
\[
i:\underline{G}_m\handwritten{\longrightarrow}G,
\]
(nécessairement central, car les composantes homogènes $V_{\underline{Q}}^n$ d'une pré-structure de Hodge sont \deleted{nécessairement}\inserted{évidemment} munis d'une \inserted{pré-}structure de Hodge de façon que $V$ soit la somme directe de $V_{\underline{Q}}^n$ en tant que pré-structure
% archive sheet number: 105
